\documentclass{article}
\usepackage[UTF8]{ctex}
\usepackage{amsmath,amssymb,amsfonts}
\usepackage{enumitem}
\usepackage{geometry}
\geometry{a4paper, margin=1in}
% === 错题本排版优化（不参与题目标识） ===
\setlength{\parindent}{0pt}
\setlength{\parskip}{0.35em plus 0.10em minus 0.08em}
\setlength{\emergencystretch}{3em}
\setlist{itemsep=0.25em, topsep=0.40em, parsep=0pt}
\allowdisplaybreaks[2]
% === 错题本排版优化结束 ===
\begin{document}
\section*{吉林大学2026年高等代数真题}
\section*{70．吉林大学2026年高等代数与解析几何试卷}

\subsection{题目1}
一．设 \(f(x)\) 是复数域上的一个多项式， \(\operatorname{deg} f(x)>0\) ，且 \(f\left(x^{2}\right) \mid f\left(x^{4}\right)\) ，证明：\(f(x)\) 的任意非零根都落在复平面的单位圆上．

\subsection{题目2}
二．设 \(A\) 为 \(n\) 阶幂等矩阵，即 \(A^{2}=A\) 。证明：

\begin{itemize}
  \item[（1）] A 可以对角化．
  \item[（2）] \(\operatorname{tr}(A)=\) 秩 \((A)\) 。
\end{itemize}

\subsection{题目3}
三。设 \(u_{1}, u_{2}, \cdots, u_{n}\) 和 \(v_{1}, v_{2}, \cdots, v_{n}\) 是内积空间 \(V\) 的两组规范正交基，\(\sigma\) 是 \(V\) 的一个等距变换，并且 \(\left[\sigma\left(u_{1}\right), \sigma\left(u_{2}\right)\right]=\left[v_{1}, v_{2}\right]\) ，即 \(\sigma\left(u_{1}\right), \sigma\left(u_{2}\right)\) 生成的子空间与 \(v_{1}, v_{2}\) 生成的子空间相等。证明：

\[
\left[\sigma\left(u_{3}\right), \sigma\left(u_{4}\right), \cdots, \sigma\left(u_{n}\right)\right]=\left[v_{3}, v_{4}, \cdots, v_{n}\right] .
\]

\subsection{题目4}
四．设 \(n\) 阶上三角矩阵

\[
A=\left(\begin{array}{ccccc}
1 & 2 & 3 & \cdots & n \\
& 1 & 2 & \cdots & n-1 \\
& & 1 & \cdots & n-2 \\
& & & \ddots & \vdots \\
& & & & 1
\end{array}\right)
\]

（1）求 \(A\) 的极小多项式和 Jordan 标准形．\\
（2）证明：\(K=\left\{B \in M_{n}(\mathrm{C}) \mid(A-I) B=O\right\}\) 是 \(M_{n}(\mathrm{C})\) 的子空间，并求 \(\operatorname{dim}_{\mathbf{C}} K\) ．

\subsection{题目5}
五．设 \(A\) 为 \(n\) 阶实对称矩阵，证明：\(A+I\) 与 \(A-I\) 合同当且仅当 \(A^{2}-I\) 正定．

\subsection{题目6}
六．在空间直角坐标系下，求平面 \(\Pi: x+y+z=2\) 内与直线 \(L:\left\{\begin{array}{l}x+2 y-2=0 \\ x+z-1=0\end{array}\right.\) 垂直相交的直线的一般方程。

\subsection{题目7}
七．在空间直角坐标系下，设 \(S_{1}\) 和 \(S_{2}\) 是如下两个球面：

\[
\begin{aligned}
& S_{1}:(x-6)^{2}+(y-8)^{2}+z^{2}=25 \\
& S_{2}:(x-10)^{2}+(y-5)^{2}+z^{2}=50
\end{aligned}
\]

求曲线 \(S_{1} \cap S_{2}\) 绕 \(z\) 轴旋转得到的曲面的方程。

\subsection{题目8}
八．在空间直角坐标系下，设双曲抛物面 \(S\) 的一般方程为 \(x^{2}-4 y^{2}=2 z\) 。

\begin{itemize}
  \item[（1）] 证明：对 \(S\) 中任意一点，有两条直母线经过它。
  \item[（2）] 求 \(S\) 中所有垂直相交的直母线的交点构成的图形的方程。
\end{itemize}
\end{document}
